% Solutions to the problems of Chapter 3 (Classical Sensor Fusion Techniques).
% All numbers are reproduced by companion_repo/exercises/ch03/solutions/ch03_hand_computations.py
% (Problems 3.1-3.9) and ch03_prob10_ekf_process_noise.py (Problem 3.10).

\begin{sol}{prob:ch03:1}
(a) With $\mathbf{G}=\begin{bmatrix}1&1\end{bmatrix}^{\mathsf T}$, $\mathbf{G}^{\mathsf T}\mathbf{G}=2$ and $\widehat{x}_{\mathrm{LS}}=\tfrac12(y_1+y_2)=42.50$\,m. The covariance formula gives $\operatorname{cov}(\widehat{x}_{\mathrm{LS}})=\tfrac14(\sigma_1^{2}+\sigma_2^{2})=\tfrac14(0.09+2.25)=0.585$\,m$^{2}$, i.e.\ a standard deviation of $0.765$\,m, worse than the radar alone.
(b) $\mathbf{\Lambda}=\sigma_1^{-2}+\sigma_2^{-2}=11.111+0.444=11.556$\,m$^{-2}$, $\widehat{x}_{\mathrm{WLS}}=\mathbf{\Lambda}^{-1}(\sigma_1^{-2}y_1+\sigma_2^{-2}y_2)=(462.22+19.29)/11.556=41.669$\,m and $\operatorname{cov}(\widehat{x}_{\mathrm{WLS}})=\mathbf{\Lambda}^{-1}=0.0865$\,m$^{2}$ ($0.294$\,m).
(c) For $\mathbf{G}=\begin{bmatrix}1&1\end{bmatrix}^{\mathsf T}$ and diagonal $\mathbf{R}$, $\mathbf{G}^{\mathsf T}\mathbf{R}^{-1}\mathbf{G}=\sigma_1^{-2}+\sigma_2^{-2}$ and $\mathbf{G}^{\mathsf T}\mathbf{R}^{-1}\mathbf{y}=\sigma_1^{-2}y_1+\sigma_2^{-2}y_2$, which is the stated average with weights $w_i=\sigma_i^{-2}/(\sigma_1^{-2}+\sigma_2^{-2})$, here $0.962$ and $0.038$. Since $\mathbf{\Lambda}\ge\sigma_1^{-2}$, $\mathbf{\Lambda}^{-1}\le\sigma_1^{2}$: an independent measurement can only add information, here reducing $0.0900$ to $0.0865$\,m$^{2}$.
(d) The ratio is $0.585/0.0865=6.8$. Ordinary least squares gives the camera the same weight as the radar although its variance is $25$ times larger; WLS gives it $3.8\,\%$ of the weight, so the fused value moves only $7$\,cm away from the radar reading.
\end{sol}

\begin{sol}{prob:ch03:2}
(a) $\det\mathbf{P}=0.16$ and $\mathbf{P}^{-1}=\begin{bmatrix}1.5625&-1.875\\-1.875&6.25\end{bmatrix}$. The measurement adds $\sigma^{-2}=11.111$ to the $(1,1)$ entry only:
$\mathbf{\Lambda}_{\mathrm{post}}=\begin{bmatrix}12.674&-1.875\\-1.875&6.25\end{bmatrix}$, $\det\mathbf{\Lambda}_{\mathrm{post}}=75.70$,
$\mathbf{P}^{+}=\begin{bmatrix}0.0826&0.0248\\0.0248&0.1674\end{bmatrix}$.
(b) $\mathbf{G}^{\mathsf T}\mathbf{R}^{-1}y=\begin{bmatrix}462.22&0\end{bmatrix}^{\mathsf T}$ and $\mathbf{P}^{-1}\mathbf{m}=\begin{bmatrix}75.0&-110.0\end{bmatrix}^{\mathsf T}$, hence $\widehat{\mathbf{x}}_{\mathrm{MAP}}=\mathbf{P}^{+}\begin{bmatrix}537.22&-110.0\end{bmatrix}^{\mathsf T}=\begin{bmatrix}41.633&-5.110\end{bmatrix}^{\mathsf T}$.
(c) $\mathbf{G}\mathbf{P}\mathbf{G}^{\mathsf T}+\mathbf{R}=1.09$, $\mathbf{K}=\begin{bmatrix}1.0&0.3\end{bmatrix}^{\mathsf T}/1.09=\begin{bmatrix}0.9174&0.2752\end{bmatrix}^{\mathsf T}$, innovation $41.6-42.0=-0.4$, so $\widehat{\mathbf{x}}=\begin{bmatrix}42.0-0.367&-5.0-0.110\end{bmatrix}^{\mathsf T}$, identical to (b); $(\mathbf{I}-\mathbf{K}\mathbf{G})\mathbf{P}$ reproduces $\mathbf{P}^{+}$.
(d) The prior correlation $0.3$ states that a range larger than predicted goes together with a less negative range rate; the second row of $\mathbf{K}$ is $P_{12}/(P_{11}+R)$, so a range innovation is transferred to the unmeasured component in proportion to the cross-covariance. With $P_{12}=0$ the gain becomes $\begin{bmatrix}0.9174&0\end{bmatrix}^{\mathsf T}$ and the range rate keeps its prior value and variance.
(e) For $\mathbf{P}^{-1}\to\mathbf{0}$ the cost reduces to the WLS cost and the estimate to the (ML) measurement; with a single range measurement the range rate would then be unobservable. For $\mathbf{R}^{-1}\to\mathbf{0}$ the measurement carries no information and $\widehat{\mathbf{x}}\to\mathbf{m}$, $\mathbf{P}^{+}\to\mathbf{P}$.
\end{sol}

\begin{sol}{prob:ch03:3}
(a) After the first measurement, $\mathbf{P}_1^{-1}=\mathbf{P}_0^{-1}+\mathbf{G}_1^{\mathsf T}\mathbf{R}_1^{-1}\mathbf{G}_1$ and $\mathbf{P}_1^{-1}\widehat{\mathbf{x}}_1=\mathbf{P}_0^{-1}\mathbf{m}+\mathbf{G}_1^{\mathsf T}\mathbf{R}_1^{-1}\mathbf{y}_1$. The second update adds $\mathbf{G}_2^{\mathsf T}\mathbf{R}_2^{-1}\mathbf{G}_2$ and $\mathbf{G}_2^{\mathsf T}\mathbf{R}_2^{-1}\mathbf{y}_2$, so that
$\mathbf{P}_2^{-1}=\mathbf{P}_0^{-1}+\mathbf{G}_1^{\mathsf T}\mathbf{R}_1^{-1}\mathbf{G}_1+\mathbf{G}_2^{\mathsf T}\mathbf{R}_2^{-1}\mathbf{G}_2$ and
$\mathbf{P}_2^{-1}\widehat{\mathbf{x}}_2=\mathbf{P}_0^{-1}\mathbf{m}+\mathbf{G}_1^{\mathsf T}\mathbf{R}_1^{-1}\mathbf{y}_1+\mathbf{G}_2^{\mathsf T}\mathbf{R}_2^{-1}\mathbf{y}_2$.
For block-diagonal $\mathbf{R}$, $\mathbf{G}^{\mathsf T}\mathbf{R}^{-1}\mathbf{G}=\sum_n\mathbf{G}_n^{\mathsf T}\mathbf{R}_n^{-1}\mathbf{G}_n$ and $\mathbf{G}^{\mathsf T}\mathbf{R}^{-1}\mathbf{y}=\sum_n\mathbf{G}_n^{\mathsf T}\mathbf{R}_n^{-1}\mathbf{y}_n$, which are exactly the batch normal equations $\mathbf{\Lambda}_{\mathrm{post}}\widehat{\mathbf{x}}_{\mathrm{ReLS}}=\mathbf{G}^{\mathsf T}\mathbf{R}^{-1}\mathbf{y}+\mathbf{P}_0^{-1}\mathbf{m}$.
(b) Both sides are sums, hence invariant under permutation of the measurements.
(c) $K_2=P_1/(P_1+\sigma_2^{2})=0.09/2.34=0.0385$, $\widehat{x}_2=41.6+0.0385\,(43.4-41.6)=41.669$\,m and $P_2=(1-K_2)P_1=0.0865$\,m$^{2}$, the WLS values of Problem~\ref{prob:ch03:1}.
(d) The additivity of information requires $\mathbf{R}$ to be block diagonal; with $\mathbf{R}_{12}\neq\mathbf{0}$ the cross terms of $\mathbf{G}^{\mathsf T}\mathbf{R}^{-1}\mathbf{G}$ do not split into a sum of per-measurement contributions, and processing the measurements as if they were independent counts the shared error twice, which is precisely the double counting that motivates covariance intersection in Section~\ref{subsec:track-fusion}.
\end{sol}

\begin{sol}{prob:ch03:4}
(a) $T^{4}/4=0.015625$, $T^{3}/2=0.0625$, $T^{2}=0.25$, so $\mathbf{Q}=0.64\begin{bmatrix}0.015625&0.0625\\0.0625&0.25\end{bmatrix}=\begin{bmatrix}0.01&0.04\\0.04&0.16\end{bmatrix}$.
(b) $\widehat{\mathbf{x}}_{k|k-1}=\begin{bmatrix}50-2&-4\end{bmatrix}^{\mathsf T}=\begin{bmatrix}48.0&-4.0\end{bmatrix}^{\mathsf T}$; $\mathbf{F}\mathbf{P}\mathbf{F}^{\mathsf T}=\begin{bmatrix}0.29&0.08\\0.08&0.16\end{bmatrix}$, so $\mathbf{P}_{k|k-1}=\begin{bmatrix}0.30&0.12\\0.12&0.32\end{bmatrix}$.
(c) $\nu_k=47.4-48.0=-0.6$\,m, $S_k=0.30+0.09=0.39$\,m$^{2}$, $\mathbf{K}_k=\begin{bmatrix}0.30&0.12\end{bmatrix}^{\mathsf T}/0.39=\begin{bmatrix}0.7692&0.3077\end{bmatrix}^{\mathsf T}$.
(d) $\widehat{\mathbf{x}}_{k|k}=\begin{bmatrix}48.0-0.4615&-4.0-0.1846\end{bmatrix}^{\mathsf T}=\begin{bmatrix}47.538&-4.185\end{bmatrix}^{\mathsf T}$ and $\mathbf{P}_{k|k}=(\mathbf{I}-\mathbf{K}_k\mathbf{H})\mathbf{P}_{k|k-1}=\begin{bmatrix}0.0692&0.0277\\0.0277&0.2831\end{bmatrix}$; the Joseph form gives the same matrix to four decimals.
(e) $d_k^{2}=0.36/0.39=0.923$, inside $[0.00098,\,5.02]$ and close to the expected value $1$: the measurement is consistent with the prediction.
(f) The prediction correlates position and velocity ($P_{12}=0.12$, from $T\,P_{22}$ and the $T^{3}/2$ term of $\mathbf{Q}$), so a position innovation carries evidence about the velocity through the second row of the gain, $P_{12}/S_k$. The position variance falls from $0.30$ to $0.30-0.30^{2}/0.39=0.069$\,m$^{2}$ because the measurement is three times more precise than the prediction; the velocity variance falls only from $0.32$ to $0.283$\,m$^{2}$/s$^{2}$, through the correlation alone.
\end{sol}

\begin{sol}{prob:ch03:5}
(a) With $d=\sqrt{p_x^{2}+p_y^{2}}$, $\partial d/\partial p_x=p_x/d$, $\partial d/\partial p_y=p_y/d$, $\partial\operatorname{atan2}(p_y,p_x)/\partial p_x=-p_y/d^{2}$ and $\partial\operatorname{atan2}(p_y,p_x)/\partial p_y=p_x/d^{2}$; the velocity columns vanish. These are the first two rows of \eqref{eq:ekf-radar-jacobian}.
(b) $d=40.112$\,m, $\mathbf{H}=\begin{bmatrix}0.9972&0.0748&0&0\\-0.00186&0.02486&0&0\end{bmatrix}$ (a finite-difference check agrees to $10^{-8}$).
(c) A bearing error $\delta\theta$ produces a position error $d\,\delta\theta$ perpendicular to the line of sight, so the lateral information decays as $1/d^{2}$ in the information matrix. With $\sigma_\theta=0.5^{\circ}=8.73\times10^{-3}$\,rad the lateral standard deviation is $60\tan\sigma_\theta=0.524$\,m at $60$\,m and $0.087$\,m at $10$\,m; this is the range-dependent scatter visible in Fig.~\ref{fig:ekf-radar-cutin}(a).
(d) Raw difference $-3.140-3.135=-6.275$\,rad; wrapped, $-6.275+2\pi=0.0082$\,rad. With $S\approx\sigma_\theta^{2}$ the NIS is $(6.275/0.00873)^{2}=5.2\times10^{5}$ for the raw difference and $(0.0082/0.00873)^{2}=0.88$ after wrapping. A $\chi^{2}_{1}$ gate at $5.02$ would reject a perfectly good detection in the first case; if it were not gated, the gain would apply a $-6.3$\,rad correction and the filter would diverge.
(e) The state and the predicted measurement are smooth quantities, and the prediction $\mathbf{g}(\widehat{\mathbf{x}}_{k|k-1})$ is already in $(-\pi,\pi]$ by construction of $\operatorname{atan2}$; the discontinuity arises only in the difference of two angles, so it is the innovation that must be mapped back to the principal interval before it enters \eqref{eq:kf_state_update} and \eqref{eq:nis}.
\end{sol}

\begin{sol}{prob:ch03:6}
(a) $\tilde p=Fp=\begin{bmatrix}1&0\end{bmatrix}^{\mathsf T}$ and $\tilde H=\begin{bmatrix}FH&\bar a\begin{bmatrix}T^{2}/2\\T\end{bmatrix}\end{bmatrix}=\begin{bmatrix}1&1.5&0.1\\0&1&0.2\end{bmatrix}$. The interval hull is obtained from the row sums of $|\tilde H|$: $x_1\in[-1.6,\,3.6]$, $x_2\in[-1.2,\,1.2]$.
(b) $\mathcal{S}=\{\mathbf{x}:|c^{\mathsf T}\mathbf{x}-1.4|\le0.3\}$, i.e.\ $1.1\le x_1\le1.7$.
(c) $\tilde H\tilde H^{\mathsf T}=\begin{bmatrix}3.26&1.52\\1.52&1.04\end{bmatrix}$, $c^{\mathsf T}\tilde H\tilde H^{\mathsf T}c=3.26$, so $\lambda^{\star}=\begin{bmatrix}3.26&1.52\end{bmatrix}^{\mathsf T}/(3.26+0.09)=\begin{bmatrix}0.9731&0.4537\end{bmatrix}^{\mathsf T}$. The residual is $y-c^{\mathsf T}\tilde p=0.4$, hence $\hat p=\begin{bmatrix}1.389&0.181\end{bmatrix}^{\mathsf T}$ and
$\hat H=\begin{bmatrix}(I-\lambda^{\star}c^{\mathsf T})\tilde H&\bar\phi\lambda^{\star}\end{bmatrix}=\begin{bmatrix}0.0269&0.0403&0.0027&0.2919\\-0.4537&0.3194&0.1546&0.1361\end{bmatrix}$,
with interval hull $x_1\in[1.028,\,1.751]$, $x_2\in[-0.882,\,1.245]$ and $\|\hat H\|_F^{2}=0.438$.
(d) With $\lambda=c$: $\hat p=\begin{bmatrix}1.4&0\end{bmatrix}^{\mathsf T}$, $\hat H=\begin{bmatrix}0&0&0&0.3\\0&1&0.2&0\end{bmatrix}$, hull $x_1\in[1.1,\,1.7]$, $x_2\in[-1.2,\,1.2]$, $\|\hat H\|_F^{2}=1.13$. The choice $\lambda=c$ reproduces the strip exactly in the measured direction but learns nothing about the velocity; $\lambda^{\star}$ minimises the Frobenius radius, which is a global size measure, and trades $0.12$ of width in $x_1$ for a reduction from $2.4$ to $2.13$ in $x_2$, using the correlation between position and velocity encoded in the generators.
(e) Maximising and minimising $x_2=\tilde H_{2,:}\xi$ over $\xi\in[-1,1]^{3}$ subject to $1.1\le1+\tilde H_{1,:}\xi\le1.7$ gives the exact hull $x_1\in[1.1,\,1.7]$, $x_2\in[-0.733,\,1.2]$. Both enclosures contain it, as they must; the $\lambda^{\star}$ enclosure is the tighter one in $x_2$ but neither is exact, because the feasible set is a polytope with slanted faces that no zonotope with these generator directions can represent.
\end{sol}

\begin{sol}{prob:ch03:7}
(a) The conflicting products are $m_C(P)m_L(V)+m_C(P)m_L(S)+m_C(\{P,V\})m_L(S)=0.20+0.05+0.03$, so $K=0.28$.
(b) The non-empty intersections collect $0.25$ on $P$ ($0.5\cdot0.4+0.5\cdot0.1$), $0.20$ on $V$ ($0.3\cdot0.4+0.2\cdot0.4$), $0.23$ on $\{P,V\}$ ($0.12+0.03+0.08$), $0.02$ on $S$ and $0.02$ on $\Omega$; their sum is $0.72=1-K$. Dividing by $0.72$: $m_D(P)=0.347$, $m_D(V)=0.278$, $m_D(\{P,V\})=0.319$, $m_D(S)=0.028$, $m_D(\Omega)=0.028$.
(c) $\mathrm{Bel}(P)=0.347$, $\mathrm{Pl}(P)=0.347+0.319+0.028=0.694$; $\mathrm{Bel}(V)=0.278$, $\mathrm{Pl}(V)=0.625$; $\mathrm{Bel}(S)=0.028$, $\mathrm{Pl}(S)=0.056$. Pignistic: $\mathrm{BetP}(P)=0.347+0.319/2+0.028/3=0.516$, $\mathrm{BetP}(V)=0.447$, $\mathrm{BetP}(S)=0.037$. The object is almost certainly a road user, but the pedestrian/vehicle decision remains open ($0.319$ on $\{P,V\}$).
(d) Smets' rule keeps the unnormalised products: $m_S(\emptyset)=0.28$, $m_S(P)=0.25$, $m_S(V)=0.20$, $m_S(\{P,V\})=0.23$, $m_S(S)=m_S(\Omega)=0.02$. The two rules use the same products; Dempster's rule redistributes the conflicting mass proportionally over the agreeing propositions, whereas Smets' rule leaves it on $\emptyset$ as an explicit measure of disagreement that a supervisor can monitor.
(e) $K=1-0.01\cdot0.01=0.9999$. Dempster's rule returns $m_D(S)=1$: two sources that are each $99\,\%$ sure of a different road user produce certainty in the one hypothesis that both consider almost impossible. Smets' rule returns $m_S(\emptyset)=0.9999$ and $m_S(S)=0.0001$, which correctly signals that the sources are irreconcilable and that at least one of them (or the frame) is unreliable; discounting or a reliability check should precede any decision.
\end{sol}

\begin{sol}{prob:ch03:8}
(a) Let $\bar{\boldsymbol{\Sigma}}$ be the joint covariance of $(\mathbf{e}_1,\mathbf{e}_2)$. For any $\mathbf{a},\mathbf{b}$ and $\omega\in(0,1)$, with $\mathbf{a}'=\sqrt{(1-\omega)/\omega}\,\mathbf{a}$ and $\mathbf{b}'=\sqrt{\omega/(1-\omega)}\,\mathbf{b}$,
$\tfrac{1-\omega}{\omega}\mathbf{a}^{\mathsf T}\bar{\mathbf{P}}_1\mathbf{a}-2\mathbf{a}^{\mathsf T}\bar{\mathbf{P}}_{12}\mathbf{b}+\tfrac{\omega}{1-\omega}\mathbf{b}^{\mathsf T}\bar{\mathbf{P}}_2\mathbf{b}=\begin{bmatrix}\mathbf{a}'\\-\mathbf{b}'\end{bmatrix}^{\mathsf T}\bar{\boldsymbol{\Sigma}}\begin{bmatrix}\mathbf{a}'\\-\mathbf{b}'\end{bmatrix}\ge0$,
which is the block inequality of the hint. Writing $\mathbf{e}=\mathbf{P}_{\mathrm{CI}}\mathbf{A}\begin{bmatrix}\mathbf{e}_1\\ \mathbf{e}_2\end{bmatrix}$ with $\mathbf{A}=\begin{bmatrix}\omega\mathbf{P}_1^{-1}&(1-\omega)\mathbf{P}_2^{-1}\end{bmatrix}$, the true covariance is
\begin{align*}
\mathbf{P}_{\mathrm{CI}}\mathbf{A}\bar{\boldsymbol{\Sigma}}\mathbf{A}^{\mathsf T}\mathbf{P}_{\mathrm{CI}}
&\preceq\mathbf{P}_{\mathrm{CI}}\left(\omega\mathbf{P}_1^{-1}\bar{\mathbf{P}}_1\mathbf{P}_1^{-1}+(1-\omega)\mathbf{P}_2^{-1}\bar{\mathbf{P}}_2\mathbf{P}_2^{-1}\right)\mathbf{P}_{\mathrm{CI}}\\
&\preceq\mathbf{P}_{\mathrm{CI}}\left(\omega\mathbf{P}_1^{-1}+(1-\omega)\mathbf{P}_2^{-1}\right)\mathbf{P}_{\mathrm{CI}}=\mathbf{P}_{\mathrm{CI}},
\end{align*}
using $\bar{\mathbf{P}}_i\preceq\mathbf{P}_i$ in the last step. The cases $\omega\in\{0,1\}$ return one of the inputs and are trivially consistent.
(b) $\mathbf{P}_{\mathrm{CI}}^{-1}(\omega)=\diag(4\omega+(1-\omega),\,\omega+6.25(1-\omega))=\diag(1+3\omega,\,6.25-5.25\omega)$, so $\tr\mathbf{P}_{\mathrm{CI}}=(1+3\omega)^{-1}+(6.25-5.25\omega)^{-1}$. Stationarity, $3(6.25-5.25\omega)^{2}=5.25(1+3\omega)^{2}$, gives $(6.25-5.25\omega)/(1+3\omega)=\sqrt{1.75}=1.3229$ and $\omega^{\star}=(6.25-1.3229)/(5.25+3\cdot1.3229)=0.5345$. Then $\mathbf{P}_{\mathrm{CI}}(\omega^{\star})=\diag(0.384,\,0.290)$ with trace $0.675$ (against $1.25$ and $1.16$ at the end points).
(c) Independent fusion gives $\mathbf{P}_{\mathrm{f}}=\diag(1/5,\,1/7.25)=\diag(0.200,\,0.138)$, trace $0.338$, half the CI value: this is the price of not knowing the correlation. CI nevertheless remains direction dependent because $\omega$ multiplies whole information matrices: the fused longitudinal variance is dominated by the radar ($\omega\cdot4$ against $(1-\omega)\cdot1$) and the lateral one by the camera ($(1-\omega)\cdot6.25$ against $\omega\cdot1$), exactly as in the independent rule, only with less total information claimed.
\end{sol}

\begin{sol}{prob:ch03:9}
(a) The graph has the edges $\{1,2\},\{2,3\},\{1,3\},\{3,4\}$; it is connected, so the priorities are identifiable up to scale. The cycle is inconsistent: $A_{12}A_{23}=6\neq A_{13}=4$.
(b) Differentiating \eqref{eq:lls-ahp} with respect to $y_i$ gives $\sum_{j\sim i}(y_i-y_j)=\sum_{j\sim i}\ln A_{ij}$, i.e.\ $\mathbf{L}\bm{y}=\bm{b}$ with
$\mathbf{L}=\begin{bmatrix}2&-1&-1&0\\-1&2&-1&0\\-1&-1&3&-1\\0&0&-1&1\end{bmatrix}$ and
$\bm{b}=\begin{bmatrix}\ln3+\ln4\\ -\ln3+\ln2\\ -\ln2-\ln4+\ln2\\ -\ln2\end{bmatrix}=\begin{bmatrix}2.485\\-0.405\\-1.386\\-0.693\end{bmatrix}$.
(c) $\mathbf{L}$ is singular ($\mathbf{L}\mathbf{1}=\mathbf{0}$); adding $\sum_iy_i=0$ gives $\bm{y}^{*}=\begin{bmatrix}1.175&0.211&-0.347&-1.040\end{bmatrix}^{\mathsf T}$ and $\bm{x}^{*}=\begin{bmatrix}0.585&0.223&0.128&0.064\end{bmatrix}^{\mathsf T}$, ranking $1\succ2\succ3\succ4$.
(d) Fitted ratios $x_1^{*}/x_2^{*}=2.62$, $x_2^{*}/x_3^{*}=1.75$, $x_1^{*}/x_3^{*}=4.58$, $x_3^{*}/x_4^{*}=2.00$; the log-residuals are $+0.135$, $+0.135$, $-0.135$ and $0$. The inconsistency $\ln(6/4)=0.405$ is split equally, $0.405/3$, over the three edges of the cycle; edge $(3,4)$ belongs to no cycle, so node $4$ is free to absorb its comparison exactly.
(e) $\widehat{A}_{14}=x_1^{*}/x_4^{*}=9.16$ and $\widehat{A}_{24}=x_2^{*}/x_4^{*}=3.49$.
(f) Without $A_{13}$ the graph is a spanning tree: the LLS fit is exact (all residuals zero), $\bm{x}^{*}\propto\begin{bmatrix}12&4&2&1\end{bmatrix}$, i.e.\ $\begin{bmatrix}0.632&0.211&0.105&0.053\end{bmatrix}$, and no consistency check of the expert is possible because every comparison is reproduced whatever its value.
\end{sol}

\begin{sol}{prob:ch03:10}
The script \repopath{ch03\_\allowbreak prob10\_\allowbreak ekf\_\allowbreak process\_\allowbreak noise.py} in the folder \repopath{exercises/\allowbreak ch03/\allowbreak solutions/} of the companion repository regenerates the scenario with a fixed seed (the realisation therefore differs slightly from the one plotted in Fig.~\ref{fig:ekf-radar-cutin}) and gives, with the constant-velocity (CV) segments and the lane change (LC) reported separately:
\begin{center}
\footnotesize
\setlength{\tabcolsep}{4pt}
\begin{tabular}{r cc cc cc}
\hline
$\sigma_a$ [m/s$^{2}$] & RMSE CV [m] & RMSE LC [m] & NIS CV & NIS LC & in $95\,\%$ CV & in $95\,\%$ LC\\
\hline
$0.1$ & $0.217$ & $0.710$ & $4.90$ & $11.54$ & $82\,\%$ & $51\,\%$\\
$0.2$ & $0.213$ & $0.538$ & $4.42$ & $6.98$ & $84\,\%$ & $70\,\%$\\
$0.4$ & $0.188$ & $0.379$ & $3.45$ & $4.30$ & $89\,\%$ & $87\,\%$\\
$0.8$ & $0.165$ & $0.260$ & $2.78$ & $3.05$ & $93\,\%$ & $95\,\%$\\
$1.6$ & $0.162$ & $0.177$ & $2.41$ & $2.36$ & $94\,\%$ & $95\,\%$\\
$3.2$ & $0.170$ & $0.133$ & $2.12$ & $1.92$ & $94\,\%$ & $93\,\%$\\
$6.4$ & $0.181$ & $0.123$ & $1.88$ & $1.66$ & $93\,\%$ & $95\,\%$\\
\hline
\end{tabular}
\end{center}
(a) For $\sigma_a\le0.4$ the filter is overconfident: the mean NIS exceeds the expected value $3$ in both segments, dramatically so during the manoeuvre ($11.5$ at $\sigma_a=0.1$, with half of the samples outside the $95\,\%$ interval), because the reported covariance does not admit the acceleration that the target actually performs. For $\sigma_a\ge1.6$ the filter is underconfident: the mean NIS falls below $3$ (down to $1.7$) because the covariance is inflated beyond what the innovations show, even though the coverage of the $95\,\%$ interval looks acceptable. The RMSE alone would not reveal either regime.
(b) On the CV segments the error is a balance between measurement noise, which is averaged out more effectively when $\mathbf{Q}$ is small, and the initialisation transient and residual model errors, which are absorbed faster when $\mathbf{Q}$ is large; the minimum ($0.162$\,m at $\sigma_a=1.6$) separates the two effects, and beyond it the filter follows the radar noise ($\sigma_y$ grows from $0.03$ to $0.07$\,m). Inside the lane change the only relevant error is the lag of a CV model behind an accelerating target, which decreases monotonically as the filter is allowed to trust the measurements more.
(c) No. The value $\sigma_a=0.8$ used in the chapter is consistent during the manoeuvre ($3.05$) but already underconfident on the CV segments ($2.78$), and the NIS peak ($12$) exceeds the interval for every $\sigma_a$: the mismatch is structural, since one $\mathbf{Q}$ is asked to describe two motion regimes. This is the argument for an interacting multiple-model filter, which switches the process noise with the detected regime, and for monitoring the NIS as a signal in its own right, as the closing remark of Example~\ref{ex:ekf-radar-cutin} anticipates.
\end{sol}

